% Part: first-order-logic
% Chapter: proof-systems
% Section: introduction

\documentclass[../../../include/open-logic-section]{subfiles}

\begin{document}

\iftag{FOL}
      {\olfileid{fol}{prf}{int}}
      {\olfileid{pl}{prf}{int}}

\olsection{تعارف}

منطقاں وچ عام طور تے معنیات وی ہُندی اے تے !!a{derivation} نظام وی۔
معنیات دا واسطہ صدق، پورا ہون جوگتا، منطقی اعتبار تے منطقی نتیجے
ورگے تصورات نال اے۔ !!{derivation} نظاماں دا مقصد منطقی نتیجہ تے
اعتبار قائم کرن دا سراسر نحوی طریقہ دینا اے۔ ایہ ایس معنی وچ سراسر
نحوی نیں کہ ایہو جہے نظام وچ کوئی !!{derivation} اک محدود نحوی شے
ہُندی اے، عام طور تے !!{sentence}s یا !!{formula}s دا کوئی سلسلہ (یا
ہور محدود بندوبست)۔ چنگے !!{derivation} نظاماں وچ ایہ خاصیت ہُندی اے
کہ !!{sentence}s یا !!{formula}s دے دِتے ہوئے ہر سلسلے یا بندوبست نوں
مشینی طریقے نال جانچیا جا سکدا اے کہ اوہ ``درست'' اے۔

پہلے درجے دی منطق لئی سب توں سادے (تے تاریخی طور تے پہلے)
!!{derivation} نظام \emph{مسلّمی} سن۔ ایہو جہے نظام وچ !!{formula}s
دا کوئی سلسلہ !!a{derivation} اُدوں گِنیا جاندا اے جے اوہدے وچ ہر
وکھرا !!{formula} یا تاں ``مسلّماں'' دے کسے مقرر سیٹ وچ شامل ہووے،
یا مقرر گنتی دے ``استنباطی قاعدیاں'' وچوں کسے اک راہیں سلسلے وچ اپنے
توں پہلاں آون والے !!{formula}s توں نکلدا ہووے---تے ایہ مشینی طریقے
نال جانچیا جا سکدا اے کہ کی !!a{formula} مسلّمہ اے تے کی اوہ استنباطی
قاعدیاں وچوں کسے اک راہیں ہور !!{formula}s توں ٹھیک طرح نکلدا اے۔
مسلّمی !!{derivation} نظاماں نوں بیان کرنا---تے ماورائے نظریاتی طور
تے سنبھالنا وی---سَوکھا اے، پر اوہناں وچ !!{derivation}s نوں پڑھنا تے
سمجھنا اوکھا اے، تے اوہناں نوں بنانا وی اوکھا اے۔

ہور !!{derivation} نظام ایس مقصد نال بنائے گئے نیں کہ !!{derivation}s
نوں بنانا سَوکھا ہووے یا مکمل ہو جان توں بعد !!{derivation}s نوں
سمجھنا سَوکھا ہووے۔ مثالاں فطری استنباط، صداقتی درخت، جنہاں نوں تابلو
ثبوت وی آکھیا جاندا اے، تے سیکوئنٹ حساب نیں۔ کجھ !!{derivation} نظام
خاص طور تے مشینی عمل نوں سامنے رکھ کے تیار کیتے گئے نیں، مثلاً
ریزولیوشن دا طریقہ سافٹ ویئر وچ نافذ کرنا سَوکھا اے (پر اوہدیاں
!!{derivation}s نوں سمجھنا لگ بھگ ناممکن اے)۔ ایہناں ہور !!{derivation}
نظاماں وچوں زیادہ تر، !!{derivation}s نوں سلسلیاں دی بجائے !!{formula}s
دے درختاں دے طور تے پیش کردے نیں۔ ایس نال ایہ ویکھنا سَوکھا ہو جاندا
اے کہ !!a{derivation} دے کیہڑے حصے ہور کیہڑے حصیاں اُتے نربھر نیں۔

سو کسی دِتّی منطق، جیویں پہلے درجے دی منطق، لئی وکھ وکھ
!!{derivation} نظام ایس گل دیاں وکھ وکھ واضح تعریفاں دیندے نیں کہ
!!a{sentence} دا \emph{قضیہ} ہون دا کی مطلب اے تے !!a{sentence} دا
کجھ ہور جملیاں توں !!{derivable} ہون دا کی مطلب اے۔ ایہ کم جیویں وی
کیتا جاوے (مسلّمی !!{derivation}s، فطری استنباطاں، سیکوئنٹ
!!{derivation}s، صداقتی درختاں یا ریزولیوشن نال ردّاں راہیں)، اسیں
چاہندے آں کہ ایہ تعلق منطقی اعتبار تے منطقی نتیجے دے معنوی تصورات نال
میل کھان۔ آؤ $\Proves !A$ نوں ``$!A$ اک قضیہ اے'' لئی تے
``$\Gamma \Proves !A$'' نوں ``$!A$،~$\Gamma$ توں !!{derivable} اے''
لئی لکھئیے۔ $\Proves$ دی تعریف جیویں وی کیتی جاوے، اسیں چاہندے آں کہ
اوہ $\Entails$ نال میل کھاوے، یعنی:
\begin{enumerate}
\item $\Proves !A$ اُدوں تے صرف اُدوں جدوں $\Entails !A$
\item $\Gamma \Proves !A$ اُدوں تے صرف اُدوں جدوں $\Gamma \Entails !A$
\end{enumerate}
اُپر دتے ``صرف اُدوں'' والے رخ نوں \emph{صحت مندی} آکھیا جاندا اے۔
!!^a{derivation} نظام صحت مند اے جے !!{derivability}، منطقی نتیجے (یا
اعتبار) دی ضمانت دیندی اے۔ ہر ڈھنگ دے !!{derivation} نظام دا صحت مند
ہونا لازم اے؛ غیر صحت مند !!{derivation} نظام بالکل فائدہ مند نہیں۔
آخرکار، !!a{derivation} دا پورا مقصد اعتبار یا منطقی نتیجے دی نحوی
ضمانت دینا اے۔ اسیں اپنے پیش کیتے !!{derivation} نظاماں دی صحت مندی
ثابت کراں گے۔

اُلٹا ``اُدوں'' والا رخ وی اہم اے: ایہنوں \emph{تمامیت} آکھیا جاندا
اے۔ تمامیت والا !!{derivation} نظام اینا طاقتور ہُندا اے کہ جدوں وی $!A$
منطقی طور تے معتبر ہووے، اوہ $!A$ نوں قضیہ وکھا سکے، تے
$\Gamma \Proves !A$ وی وکھا سکے جدوں وی $\Gamma \Entails !A$۔ تمامیت
قائم کرنا ودھ اوکھا اے، تے کجھ منطقاں دا کوئی تمامیت والا !!{derivation} نظام
نہیں۔ پہلے درجے دی منطق دا اے۔ کرٹ گوڈل اوہ پہلا شخص سی جہنے 1929 دے
اپنے ڈاکٹری مقالے وچ پہلے درجے دی منطق دے !!a{derivation} نظام لئی
تمامیت ثابت کیتی۔

!!{derivation} نظاماں نال جڑیا اک ہور تصور \emph{بے تضادی} دا اے۔
!!{sentence}s دے سیٹ نوں متضاد آکھیا جاندا اے جے اوس توں جو کجھ وی
!!{derive} کیتا جا سکدا ہووے، تے نہیں تاں بے تضاد۔ متضاد ہونا، پورا
ہون جوگ نہ ہون دا نحوی ہم منصب اے: پورا ہون جوگ نہ ہون والے سیٹاں
وانگ، !!{sentence}s دے متضاد سیٹ چنگے نظریے نہیں بناؤندے؛ اوہ بنیادی
طور تے خراب ہُندے نیں۔ !!{sentence}s دے بے تضاد سیٹ ضروری نہیں کہ سچے
یا فائدہ مند ہون، پر گھٹّو گھٹ منطقی فائدے دی ایہ سب توں مُڈھلی حد
پار کردے نیں۔ وکھ وکھ !!{derivation} نظاماں لئی !!{sentence}s دے سیٹ
دی بے تضادی دی خاص تعریف وکھری ہو سکدی اے، پر~$\Proves$ وانگ اسیں
چاہندے آں کہ بے تضادی اپنے معنوی ہم منصب، پورا ہون جوگتا، نال اکّو
ہووے۔ اسیں چاہندے آں کہ ہمیشہ ایہ گل ہووے کہ $\Gamma$ اُدوں تے صرف
اُدوں بے تضاد اے جدوں اوہ پورا ہون جوگ اے۔ ایتھے ``صرف اُدوں'' والا
رخ تمامیت دے برابر اے (بے تضادی، پورا ہون جوگتا دی ضمانت دیندی اے)،
تے ``اُدوں'' والا رخ صحت مندی دے برابر اے (پورا ہون جوگتا، بے تضادی
دی ضمانت دیندی اے)۔ اصل وچ، کلاسیکی پہلے درجے دی منطق لئی صحت مندی
تے تمامیت دے دونیں روپ ہم ارز نیں۔

\end{document}
